\lecture{Лекция 8}{}

\subsection{Метод отсекающей гиперплоскости}
\label{subsec:cutting-plane}

\begin{idea}[title={Обозначения}]
В этом параграфе $\Omega$ — допустимое множество \emph{исходной} задачи,
а $S_k$ — аппроксимирующие его многогранники. В тетради оба обозначены одной
буквой, что мешает читать выкладки.
\end{idea}

Рассматривается задача
\[
  \begin{cases}
    \varphi_0(x) \to \min , & x \in \Omega , \\
    \Omega = \{\, x \in \R^n \mid \varphi(x) \le 0 \,\} ,
  \end{cases}
\]
где $\varphi$ — выпуклая дифференцируемая, а целевая функция линейна:
$\varphi_0(x) = c\T x$, $c \in \R^n$.

\subsubsection*{Приведение задачи к этому виду}

\begin{enumerate}[label=\arabic*),left=0pt]
  \item Если допустимое множество задано системой
        $\Omega = \{\, x \mid \varphi_i(x) \le 0,\ i = \overline{1,m} \,\}$,
        то достаточно взять
        \[
          \varphi(x) = \max_{i = \overline{1,m}} \varphi_i(x)
        \]
        — одно ограничение вместо $m$.
  \item Если $\varphi_0(x)$ не является линейной (но выпукла), вводим
        вспомогательную переменную $x^{(n+1)}$ и задачу
        \[
          \begin{cases}
            x^{(n+1)} \to \min , \\
            \varphi(x) \le 0 , \\
            \varphi_0(x) - x^{(n+1)} \le 0 .
          \end{cases}
        \]
  \item Объединяя 1) и 2) — сначала реализуется 2), потом 1):
        \[
          \begin{cases}
            x^{(n+1)} \to \min , \\
            \varphi_i\big(x, x^{(n+1)}\big) \le 0 , \quad i = \overline{1,m} , \\
            \varphi_{m+1}\big(x, x^{(n+1)}\big) = \varphi_0(x) - x^{(n+1)} ,
          \end{cases}
          \qquad
          \varphi(x) = \max_{i} \varphi_i\big(x, x^{(n+1)}\big) ,
        \]
        то есть
        \[
          \begin{cases}
            x^{(n+1)} \to \min , \\
            \varphi\big(x, x^{(n+1)}\big) \le 0 .
          \end{cases}
        \]
\end{enumerate}

\begin{idea}
Зачем эти преобразования. Весь метод построен на том, что на каждой итерации
решается \emph{линейная} задача на многограннике — иначе никакого выигрыша от
аппроксимации нет. Пункт~2) — стандартный приём «эпиграфной» переформулировки:
минимизировать выпуклую $\varphi_0$ — то же самое, что минимизировать новую
координату $x^{(n+1)}$ при условии $\varphi_0(x) \le x^{(n+1)}$. Пункт~1)
сворачивает систему неравенств в одно, пользуясь тем, что максимум конечного
числа выпуклых функций выпуклен. После обоих шагов задача имеет ровно тот вид,
для которого описан алгоритм.
\end{idea}

\subsubsection*{Алгоритм}

Считаем, что $\Omega \ne \varnothing$ и $\Omega$ — компакт.

\begin{algo}[title={Метод отсекающей гиперплоскости}]
\textbf{I. Начальный этап.} Построим объемлющий многогранник
\[
  S = \big\{\, x \mid a_i\T x - b^{(i)} \le 0,\ i = \overline{-1,0} \,\big\} ,
  \qquad S \ \text{— компакт}, \quad \Omega \subset S .
\]
Положим $S_0 = S$, $k = 0$, решим вспомогательную задачу
$\varphi_0 \to \min$, $x \in S_0$; пусть $x_0$ — её решение.

\medskip
\textbf{II. Основной этап ($k$-й шаг).} Пусть $S_k$ построено; решаем
$\varphi_0(x) \to \min$, $x \in S_k$, и получаем $x_k$.

Строим \textbf{субградиент} $a_{k+1}$ функции $\varphi$ в точке $x_k$:
\[
  \varphi(x) \ \ge\ \varphi(x_k) + a_{k+1}\T (x - x_k) \qquad \forall x .
\]
Обозначив $b^{(k+1)} = -\varphi(x_k) + a_{k+1}\T x_k$, получаем
\[
  \varphi(x) \ \ge\ a_{k+1}\T x - b^{(k+1)} ,
\]
и определяем
\[
  S_{k+1} = S_k \cap \big\{\, a_{k+1}\T x - b^{(k+1)} \le 0 \,\big\}
          = \big\{\, x \in S_k \ \big|\ a_{k+1}\T x - b^{(k+1)} \le 0 \,\big\} .
\]

\medskip
\textbf{III. Итерационный переход:} $k \to k+1$.

\medskip
\textbf{IV. Условие останова:} $\norm{x_{k+1} - x_k} < \eps$.
\end{algo}

\inputfig{figures/fig-cutting-subgradient}

\begin{idea}
Почему отсечение корректно, то есть не задевает $\Omega$. Для любой допустимой
точки $x \in \Omega$ выполнено $\varphi(x) \le 0$, а неравенство субградиента
даёт $a_{k+1}\T x - b^{(k+1)} \le \varphi(x) \le 0$ — новое ограничение
выполняется автоматически. Текущая же точка $x_k$ отсекается: для неё
$a_{k+1}\T x_k - b^{(k+1)} = \varphi(x_k) > 0$, если $x_k \notin \Omega$. То
есть каждая итерация выбрасывает заведомо недопустимую точку, не теряя ни одной
допустимой. Это и доказывается ниже в лемме~\ref{lem:cut-nested}.
\end{idea}

\subsubsection*{Особенности алгоритма}

\paragraph{I. Линейная подзадача.}
На каждом итерационном шаге решается
\[
  \begin{cases}
    c\T x \to \min , \\
    a_i\T x - b^{(i)} \le 0 , \quad i = \overline{-1, k} ,
  \end{cases}
  \tag{$*$}\label{eq:cut-lp}
\]
то есть от шага к шагу меняется \emph{только число ограничений} — значит
применим симплекс-метод (\S4).

\paragraph{II. Двойственная задача.}
Задача~\eqref{eq:cut-lp} равносильна
\[
  \begin{cases}
    c\T x \to \min , \\
    -a_i\T x \ge -b^{(i)} , \quad i = \overline{-1,k} ,
  \end{cases}
\]
двойственная к которой —
\[
  \begin{cases}
    \displaystyle -\sum_{i = -1}^{k} b^{(i)} y^{(i)} \ \to\ \max
      \qquad\text{то есть}\qquad -b\T y \to \max , \\[8pt]
    \displaystyle -\sum_{i = -1}^{k} y^{(i)} a_i = c ,
    \qquad y^{(i)} \ge 0 , \quad i = \overline{-1,k} .
  \end{cases}
\]

\begin{idea}
Здесь и прячется главный вычислительный выигрыш. Добавление одного ограничения
к прямой задаче — это добавление одной \emph{переменной} к двойственной.
Поэтому предыдущее двойственное решение, дополненное нулевой компонентой,
остаётся допустимым для новой двойственной задачи и служит готовым начальным
приближением: симплекс-метод продолжает работу с того места, где остановился,
а не стартует заново. Это и есть смысл фразы «в качестве начального приближения
даёт $y$ текущей итерации, дополненное нулевой компонентой».
\end{idea}

\paragraph{III. Построение субградиента.}
Для построения $S_{k+1}$ ищется субградиент $\varphi(x)$ в точке $x_k$:
\[
  \varphi(x) \ \ge\ \varphi(x_k) + a_{k+1}\T (x - x_k) .
\]
Возможны два случая.

\begin{enumerate}[label=\arabic*),left=0pt]
  \item Задача пришла в исходной постановке с \emph{дифференцируемой}
        $\varphi(x)$ — тогда градиент совпадает с субградиентом,
        $a_{k+1} = \nabla\varphi(x_k)$.
  \item Задача сведена к максимуму,
        $\varphi(x) = \max\limits_{i=\overline{1,m}} \varphi_i(x)$. Обозначим
        \[
          I(x_k) = \Big\{\, i \ \Big|\ \varphi_i(x_k)
            = \max_{i=\overline{1,m}} \varphi_i(x_k) \,\Big\}
        \]
        — набор индексов, на которых максимум достигается. По этим ограничениям
        строится субградиент (все $\varphi_i$ дифференцируемы):
        \begin{equation}\label{eq:subgrad-max}
          a_{k+1} = \sum_{i \in I(x_k)} \lambda_i \nabla\varphi_i(x_k) ,
          \qquad
          \lambda_i \ge 0 \ \ \forall i \in I(x_k) ,
          \qquad
          \sum_{i \in I(x_k)} \lambda_i = 1 .
        \end{equation}
\end{enumerate}

\begin{lemma}\label{lem:subgrad-max}
Пусть $\varphi_i$, $i = \overline{1,m}$, выпуклы и дифференцируемы,
$\varphi(x) = \max_i \varphi_i(x)$, и
$I(x) = \{\, i \mid \varphi_i(x) = \varphi(x) \,\}$. Тогда всякий вектор вида
\[
  a = \sum_{i \in I(x)} \lambda_i \nabla\varphi_i(x) ,
  \qquad \lambda_i \ge 0 , \qquad \sum_{i \in I(x)} \lambda_i = 1 ,
\]
является субградиентом $\varphi$ в точке $x$.
\end{lemma}

\begin{idea}[title={Доказательство (на лекции не приводилось)}]
\begin{proof}
Возьмём произвольную точку $z \in \R^n$ и любой индекс $j \in I(x)$. Так как
$\varphi$ — максимум, а $\varphi_j$ выпукла и дифференцируема,
\[
  \varphi(z) \ \ge\ \varphi_j(z)
  \ \ge\ \varphi_j(x) + \nabla\T\varphi_j(x)\,(z - x) .
\]
Но $j \in I(x)$ означает $\varphi_j(x) = \varphi(x)$, поэтому
\[
  \varphi(z) \ \ge\ \varphi(x) + \nabla\T\varphi_j(x)\,(z - x)
  \qquad \forall j \in I(x) .
\]
Умножим неравенство с номером $j$ на $\lambda_j \ge 0$ и просуммируем по
$j \in I(x)$, пользуясь $\sum_j \lambda_j = 1$:
\[
  \varphi(z) = \sum_{j \in I(x)} \lambda_j\, \varphi(z)
  \ \ge\ \sum_{j \in I(x)} \lambda_j
         \Big( \varphi(x) + \nabla\T\varphi_j(x)(z-x) \Big)
  = \varphi(x) + a\T (z - x) .
\]
Это и есть определение субградиента. \qedhere
\end{proof}

\medskip
Обратное утверждение (всякий субградиент максимума имеет вид~\eqref{eq:subgrad-max},
то есть $\partial\varphi(x) = \operatorname{conv}\{\nabla\varphi_i(x)\}_{i \in I(x)}$)
тоже верно, но его доказательство существенно длиннее и на лекции не нужно: для
алгоритма достаточно уметь \emph{предъявить} хотя бы один субградиент, а проще
всего взять $\lambda_{i_0} = 1$ для какого-нибудь одного $i_0 \in I(x_k)$, то
есть $a_{k+1} = \nabla\varphi_{i_0}(x_k)$.
\end{idea}

\subsubsection*{Сходимость}

\begin{lemma}\label{lem:cut-nested}
$\Omega \subset S_k$ для всех $k \ge 1$.
\end{lemma}

\begin{proof}
Возьмём произвольный элемент $x \in \Omega$ и покажем, что он лежит в $S_k$.
По определению $x \in \Omega \iff \varphi(x) \le 0$. С другой стороны, по
неравенству субградиента для каждого $j = \overline{1,k}$
\[
  \underbrace{\varphi(x)}_{\le\ 0}
  \ \ge\ \varphi(x_{j-1}) + a_j\T (x - x_{j-1})
  = a_j\T x - b^{(j)} ,
\]
откуда $a_j\T x - b^{(j)} \le 0$ для всех $j = \overline{1,k}$, то есть
$x \in S_k$.
\end{proof}

\begin{remark}
Целевая функция вдоль итераций \emph{не убывает}:
\[
  \varphi_0(x_0) \le \varphi_0(x_1) \le \dots \le \varphi_0(x_k) \le \dots ,
\]
поскольку $S_{k+1} \subset S_k$, а минимум по меньшему множеству не меньше. При
этом по лемме~\ref{lem:cut-nested} $\Omega \subset S_k$, значит
\[
  \varphi_0(x_*) \ \ge\ \dots \ \ge\ \varphi_0(x_k) \ \ge\ \dots ,
\]
где $x_*$ — оптимальное решение исходной задачи. То есть $\varphi_0(x_k)$ —
возрастающая последовательность \emph{нижних} оценок оптимума.
\end{remark}

\begin{theorem}\label{thm:cutting-plane}
Пусть $\varphi(x)$ непрерывна и выпукла, $\Omega$ — компакт, исходная
аппроксимация $S$ — тоже компакт, и субградиенты ограничены:
\[
  \norm{a(x)} \le R , \quad x \in S , \qquad \exists R > 0 .
\]
Тогда любая предельная точка последовательности $\{x_k\}$ является оптимальной,
причём $\varphi(x_k) \to 0$.
\end{theorem}

\begin{proof}
\textbf{1) Предельные точки существуют.} $S_0 = S$ и $S_{k+1} \subset S_k$,
значит вся последовательность $\{x_k\}$ лежит в компакте $S$, и у неё есть
сходящиеся подпоследовательности.

\medskip
\textbf{2) Если ограничение выполнено на каком-то шаге.} Пусть
$\exists k \in \NN$: $\varphi(x_k) \le 0$. Тогда $x_k \in \Omega$, и потому
$\varphi_0(x_k) \ge \varphi_0(x_*)$. С другой стороны, по
лемме~\ref{lem:cut-nested} $\Omega \subset S_k$, а $x_k$ минимизирует
$\varphi_0$ на $S_k$, поэтому
\[
  \varphi_0(x_k) = \min_{x \in S_k} \varphi_0(x)
  \ \le\ \min_{x \in \Omega} \varphi_0(x) = \varphi_0(x_*) .
\]
Следовательно $\varphi_0(x_k) = \varphi_0(x_*)$, то есть $x_k$ — оптимальное
решение, и алгоритм закончил работу.

\medskip
\textbf{3) Общий случай: $\varphi(x_k) > 0$ для всех $k$.} Покажем, что тогда
существует подпоследовательность с $\varphi(x_{k_j}) \to 0$.

От противного: пусть существуют $\beta > 0$ и бесконечное множество индексов
$J$ такие, что $\varphi(x_k) \ge \beta$ при $k \in J$. Последовательность
$\{x_k\}_{k \in J}$ лежит в компакте $S$, значит содержит фундаментальную
подпоследовательность: для любого $\delta > 0$ найдутся индексы $p < q$ из $J$ с
$\norm{x_p - x_q} \le \delta$.

Возьмём $p < q$. Так как $x_q \in S_q \subset S_{p+1}$, точка $x_q$
удовлетворяет $(p{+}1)$-му отсечению:
\[
  0 \ \ge\ a_{p+1}\T x_q - b^{(p+1)}
  = a_{p+1}\T x_q + \varphi(x_p) - a_{p+1}\T x_p
  = \varphi(x_p) + a_{p+1}\T (x_q - x_p) ,
\]
откуда
\[
  \varphi(x_p) \ \le\ a_{p+1}\T (x_p - x_q)
  \ \le\ \norm{a_{p+1}} \cdot \norm{x_p - x_q}
  \ \le\ R\,\norm{x_p - x_q} .
\]
Выберем теперь $p < q$ из $J$ так, чтобы $\norm{x_p - x_q} \le \beta/(2R)$.
Получаем $\varphi(x_p) \le \beta/2 < \beta$ — противоречие с $p \in J$.

Значит предположение неверно, и существует подпоследовательность $\{x_{k_j}\}$,
вдоль которой $\varphi(x_{k_j}) \to 0$.

\medskip
\textbf{4) Предельная точка оптимальна.} Переходя при необходимости к
подпоследовательности, считаем $x_{k_j} \to \bar{x}$. По непрерывности
$\varphi$
\[
  \lim_{j \to \infty} \varphi(x_{k_j}) = \varphi(\bar{x}) = 0
  \qquad\Longrightarrow\qquad \bar{x} \in \Omega
  \qquad\Longrightarrow\qquad \varphi_0(\bar{x}) \ge \varphi_0(x_*) .
\]
С другой стороны, как в пункте~2), $\varphi_0(x_{k_j}) \le \varphi_0(x_*)$ для
всех $j$, и по непрерывности $\varphi_0(\bar{x}) \le \varphi_0(x_*)$. Сравнивая
два неравенства,
\[
  \varphi_0(\bar{x}) = \varphi_0(x_*) ,
\]
то есть $\bar{x}$ — оптимальное решение исходной задачи.
\end{proof}

\begin{idea}
Две оценки на $\varphi_0(\bar{x})$ приходят из разных мест, и именно это стоит
держать в голове:
\begin{itemize}[nosep,left=0pt]
  \item $\varphi_0(\bar{x}) \ge \varphi_0(x_*)$ — потому что $\bar{x}$
        \emph{допустима} ($\varphi(\bar x) = 0$), а $x_*$ — минимум по всем
        допустимым точкам;
  \item $\varphi_0(\bar{x}) \le \varphi_0(x_*)$ — потому что каждая $x_{k_j}$
        минимизирует $\varphi_0$ по \emph{более широкому} множеству
        $S_{k_j} \supset \Omega$.
\end{itemize}
Метод зажимает оптимум с двух сторон: снизу — значениями на релаксациях,
сверху — допустимостью предельной точки.
\end{idea}
